\section{Introduction}
\label{sec:introduction}

Large language models (LLMs) pre-trained on vast corpora have become the dominant
paradigm for natural language processing \cite{Vaswani2017, Devlin2019}. A single
pre-trained transformer backbone can be adapted to dozens of downstream tasks through
fine-tuning, and parameter-efficient fine-tuning (PEFT) methods have reduced the cost
of this adaptation to a fraction of the original model size \cite{Houlsby2019, Hu2022}.
The result is an ecosystem in which a single frozen backbone serves as a universal
substrate, and small task-specific adapter modules provide the task-particular
knowledge required for each application.

This ecosystem creates a substantial opportunity for collaborative knowledge sharing.

If adapters are small, self-contained, and additive to a shared frozen backbone, they
are in principle \textit{transmissible}: a model trained by one party on a specialist
task can be shared with another party who needs that capability, without revealing
the underlying training data, and without requiring the receiving party to retrain the
backbone from scratch. In a peer-to-peer (P2P) network of nodes sharing the same
frozen backbone, adapters are, in principle, transmissible units of knowledge exchange. Each peer
may contribute specialist adapters it has trained, and benefit from adapters
contributed by others.

\textbf{The research question} addressed by this review is: \textit{is a
peer-to-peer research direction for adapter-based multi-task inference
theoretically grounded and supported by the current state of the literature?}
More specifically: given a network of autonomous nodes each holding a frozen
transformer backbone and a local collection of task-specific adapters, does the
existing literature provide the building blocks required for any node to
discover, acquire, compose, and serve adapters from across the network for
arbitrary multi-task NLP queries, without any central coordinator, and if
not, what are the open gaps?
Section~\ref{sec:methodology} operationalises this feasibility question into a literature-searchable form: \textit{what methods have been proposed 
for decentralised, adapter-based multi-task inference over a shared frozen 
transformer backbone, and what are the open research gaps?}

This problem sits at the intersection of three active research fields, which this
review decomposes into five sub-questions, each addressed by one chapter:
\begin{enumerate}
  \item[\textbf{SQ1}] (Section~\ref{sec:peft}): \textit{What structural properties of
    parameter-efficient adapters make them suitable as transmissible, composable
    knowledge units in a peer-to-peer setting?}
  \item[\textbf{SQ2}] (Section~\ref{sec:adapter_composition}): \textit{Can
    independently trained adapters be composed into multi-task capability without
    joint retraining or centralised coordination?}
  \item[\textbf{SQ3}] (Section~\ref{sec:inference_systems}): \textit{What serving and
    memory-management mechanisms exist for multi-adapter inference, and what
    centralisation assumptions do they carry?}
  \item[\textbf{SQ4}] (Section~\ref{sec:moe_routing}): \textit{Can routing mechanisms
    select the right adapter for a query without a fully enumerated, centrally known
    expert pool?}
  \item[\textbf{SQ5}] (Section~\ref{sec:p2p_federated}): \textit{What decentralised
    infrastructure and protocols exist for adapter-level exchange without a central
    coordinating server?}
\end{enumerate}
Section~\ref{sec:synthesis_gap} integrates the five partial answers into a unified gap
map.

The multi-tenant serving literature (S-LoRA \cite{Sheng2024}, Punica \cite{Chen2023Punica}, 
CaraServe \cite{Li2024CaraServe}) assumes a single administrative domain. 
The federated PEFT literature assumes a central server for adapter aggregation. 
The P2P inference literature (Petals \cite{Borzunov2023}) distributes backbone 
layers across nodes but does not address per-task adapter modules. The closest precedent is 
{\v{S}}ajina et al.~\cite{Sajina2024}, which demonstrates that task-specific modules can be trained and served 
across a P2P network of encoder-only transformers — but the modules are tied to a fixed peer 
population, are not framed as PEFT/LoRA-style adapters, and are neither discoverable nor composable 
by peers outside the original training group. Each of these bodies of work addresses part
of the problem, but their intersection has not yet been addressed as an integrated whole.

Prior surveys that cover these areas do not close this gap either. Comprehensive reviews of
parameter-efficient fine-tuning \cite{Han2024,WangAIReview2024} and of LoRA specifically
\cite{MaoLoraSurvey2024} treat adapters as artefacts to be trained more efficiently within a
single controlled environment; they do not frame an adapter as a transferable, exchangeable
unit of capability in a decentralised system. Conversely, surveys and position papers on
decentralised federated learning \cite{YeDecentralizedFL2022,WinkP2PFL2021} and on autonomous
P2P connection establishment \cite{SajinaP2PConnection2024} address the network and protocol
layer but remain agnostic to the semantics of the modules being exchanged. In short, the
existing reviews either assume a monolithic, centrally administered model or a module-agnostic
P2P substrate; none examines the two jointly. The specific synthesis gap this review closes is
therefore the cross-cutting question of \emph{how adapter-level knowledge can be represented,
discovered, and exchanged in a decentralised, multi-task P2P setting} --- a gap that no prior
survey or review, within the 123-paper corpus examined in
Section~\ref{sec:methodology}, addresses directly.

\subsection{Review Structure}
\label{sec:introduction:structure}

Section~\ref{sec:background} provides the technical background on transformer
architectures, PEFT fundamentals, and P2P network primitives. Section~\ref{sec:methodology}
describes the systematic literature review methodology. Sections~\ref{sec:peft}--\ref{sec:p2p_federated}
review the five bodies of related work. Section~\ref{sec:synthesis_gap} synthesises the
literature into a unified gap analysis and identifies open research directions.
Section~\ref{sec:conclusion} concludes.
